\textbf{部分递归函数}是指那些能够从基本函数 $\Zero$、$\Succ$、$\Proj{i}{j}$ 出发，通过合成、原始递归和无界搜索定义的函数。这类定义本身也可以编码为自然数。在选定一个合适的编码后，若一个部分递归函数~$f$ 的代码是~$e$，我们就称 $e$ 为~$f$ 的一个\textbf{指标}，并将 $f$ 也记作~$\cfind{e}$。不仅函数的定义，函数应用于自变元的计算过程也能编码为自然数。判定 $s$ 是否编码了 $\cfind{e}$ 应用于自变元 $x$ 的计算过程的性质 $T(e, x, s)$，以及以计算代码 $s$ 为输入、返回该计算输出的函数 $U(s)$，二者都是原始递归的。因此，我们有 $\cfind{e}(x) \simeq
U(\umin{s}{T(e,x,s)})$。这就是\textbf{Kleene（克林）范式定理}。

Kleene范式定理实际上不仅适用于部分递归函数，也适用于任何已知的计算模型。例如，对于Turing机，同样可以将其定义和计算过程编码为自然数，使得Kleene范式定理成立。因此，通过Kleene范式定理对部分可计算函数进行指标化，便能对可计算函数展开非常一般性的研究。这个研究领域被称为\textbf{可计算性理论}。其成果包括：存在一个二元通用部分递归函数~$\fn{Un}$，使得 $\fn{Un}(e,x) \simeq \cfind{e}(x)$；不存在全的通用可计算函数；以及判定 $\cfind{e}(x)$ 是否有定义的问题是不可判定的（即\textbf{停机问题}）。\textbf{自停机集}~$K = \Setabs{e}{\cfind{e}(e)
\text{ 有定义}}$ 是另一个不可判定的自然数集合的例子。然而，它是\textbf{可计算枚举的}，也就是说，它是某个可计算函数的值域。

通过Gödel编码，可计算性理论可对Gödel关于理论不可判定性与不完全性的定理给出非常一般的表述。利用可计算性理论的结果，我们可以证明：每个表示所有原始递归关系的 $\omega$-一致理论都是不可判定的，并且若这样的理论是可计算枚举的，则它必然是不完全的。
